\documentclass[12pt,a4paper]{report}

\usepackage[utf8]{inputenc}
\usepackage[T1]{fontenc}
\usepackage[spanish,es-tabla]{babel}
\usepackage{geometry}
\geometry{margin=2.5cm}
\usepackage{graphicx}
\usepackage{grffile}
\usepackage{hyperref}
\usepackage{fancyhdr}
\usepackage{booktabs}
\usepackage{longtable}
\usepackage{array}
\usepackage{multirow}
\usepackage{listings}
\usepackage{xcolor}
\usepackage{caption}
\usepackage{subcaption}
\usepackage{tikz}
\usepackage{fontspec}
\setmainfont{Arial}
\usetikzlibrary{positioning,arrows.meta,shapes.geometric,fit}
\usepackage{titlesec}
\titleformat{\section}{\large\bfseries}{\thesection}{1em}{}
\titleformat{\subsection}{\normalsize\bfseries}{\thesubsection}{1em}{}

\hypersetup{colorlinks=true, linkcolor=black, urlcolor=blue, citecolor=black}
\pagestyle{fancy}
\fancyhf{}
\fancyhead[L]{\leftmark}
\fancyhead[R]{\thepage}
\renewcommand{\chaptermark}[1]{\markboth{#1}{}}

\renewcommand{\thechapter}{\Roman{chapter}}
\renewcommand{\thesection}{\thechapter.\arabic{section}}
\renewcommand{\thesubsection}{\thesection.\arabic{subsection}}

\lstset{
  basicstyle=\ttfamily\small,
  breaklines=true,
  frame=single,
  backgroundcolor=\color{gray!10},
  keywordstyle=\color{blue},
  commentstyle=\color{gray},
  stringstyle=\color{teal},
  extendedchars=true,
  inputencoding=utf8,
  literate=
    {á}{{\'a}}1 {é}{{\'e}}1 {í}{{\'i}}1 {ó}{{\'o}}1 {ú}{{\'u}}1
    {ñ}{{\~n}}1 {Á}{{\'A}}1 {É}{{\'E}}1 {Í}{{\'I}}1 {Ó}{{\'O}}1
    {Ú}{{\'U}}1 {Ñ}{{\~N}}1,
}

\graphicspath{{imagenes/}}

\newcommand{\screenshot}[2]{%
  \begin{figure}[htbp]
    \centering
    \includegraphics[width=0.92\textwidth]{#1}
    \caption{#2}
  \end{figure}%
}

\newcommand{\screenshotsmall}[2]{%
  \begin{figure}[htbp]
    \centering
    \includegraphics[width=0.5\textwidth]{#1}
    \caption{#2}
  \end{figure}%
}

\begin{document}

% --- PORTADA ---
\begin{titlepage}
  \centering
  \vspace*{2cm}
  {\Large UNIVERSIDAD DEL BÍO-BÍO\par}
  \vspace{0.5cm}
  {\large Ingeniería de Software\par}
  \vspace{3cm}
  {\huge\bfseries Sistema de Gestión para Empresa de Aseos\par}
  \vspace{0.5cm}
  {\Large Región del Bío-Bío\par}
  \vspace{2cm}
  {\large Proyecto de Ingeniería de Software\par}
  \vspace{3cm}
\begin{center}
\begin{tabular}{r l}
  Alumno: & \textbf{Ricardo Alejandro Yáñez Ríos} \\
  
  Profesora: & \textbf{Alejandra Segura}
\end{tabular}
\end{center}
  \vfill
  {\large Julio 2026\par}
  Concepción -- Chile
\end{titlepage}

\pagenumbering{roman}
\tableofcontents
\listoffigures
\listoftables
\clearpage
\pagenumbering{arabic}

% =============================================================================
\clearpage
\chapter{Introducción}
% =============================================================================

El objetivo de este trabajo es desarrollar y documentar, el proyecto \textit{Sistema de Gestión para Empresa de Aseos}, desarrollado como parte del curso de Ingeniería de Software (ISW) buscando enfrentar estas tecnologías a los retos para una transformación digital. Este trabajo en si, busca apoyar la gestión operativa y contractual de una empresa de servicios de aseo ubicada en Concepción, cuyos clientes se concentran en la Región del Bío-Bío.

En el documento se presenta el contexto organizacional, la definición del software propuesto, el análisis y diseño de la solución, la implementación técnica realizada en el repositorio, y las conclusiones del trabajo. El desarrollo combina un backend en Node.js con Express y PostgreSQL, y un frontend en React con Vite, comunicados mediante una API REST.

\section{Presentación del contexto o área de su problemática}

\subsection{Presentación Rubro Institución/Empresa}

La empresa objeto de estudio pertenece al rubro de servicios de aseo y limpieza profesional. Opera exclusivamente en la Región del Bío-Bío, con oficina central en Concepción. Su cartera de clientes incluye instituciones públicas y privadas como por ejemplo; Hospitales, Bancos, Universidades y otras entidades que contratan servicios a través del mercado público o mediante comunicación directa.

El modelo de negocio consiste en proveer personal de aseo junto con herramientas básicas de limpieza. La cantidad de trabajadores y el tipo de implementación varían según el cliente: un Hospital requiere mayor volumen de personal y equipamiento especializado (trajes de bioseguridad, desinfectantes hospitalarios) que una sucursal bancaria, donde basta con aspiradoras y productos químicos estándar.

La empresa mantiene relaciones de largo plazo con clientes frecuentes, aunque no exclusivos. Los contratos tienen duraciones entre 1 y 5 años, con posibilidad de incorporar personal adicional capacitado en aseo cuando surgen nuevas necesidades por ejemplo; un nuevo edificio universitario o redistribuir el personal existente en cambios menores.

El personal fijo incluye: un administrador de empresas, un contador auditor, un químico, un prevencionista de riesgos, un encargado de RRHH, un encargado de ventas y personal capacitado en aseo. La administración de la empresa se encarga directamente de la organización de turnos y de la asignación de personal a cada ubicación de servicio.

\subsection{Descripción de problemas u oportunidades de mejora}

El levantamiento de información reveló las siguientes problemáticas y oportunidades de mejora:

\begin{itemize}
  \item \textbf{Gestión manual de asignaciones:} La determinación del volumen de personal y su distribución por ubicación se realiza de forma informal, dificultando el seguimiento cuando hay múltiples clientes activos.
  \item \textbf{Variabilidad por tipo de cliente:} Cada institución exige una valoración distinta o cálculos en los costos diferente por los trabajos realizados, herramientas y cantidad de personal distintos, lo que complica la planificación sin un registro centralizado.
  \item \textbf{Contratos de largo plazo:} Los acuerdos de 1 a 5 años requieren control de vigencia, ajustes por IPC y renovaciones, tareas actualmente dispersas en planillas.
  \item \textbf{Redistribución ante cambios:} Las negociaciones contractuales pueden implicar incorporar o reasignar personal, al no poseer un sistema de software, estos cambios no quedan trazados.
  \item \textbf{Contratación ejecutiva:} Posiciones críticas (administrador, contador, prevencionista, quimico, ventas u otros) necesitan un proceso de seguimiento que hoy no está sistematizado.
  \item \textbf{Presupuesto y negociación:} Cualquier desviación presupuestaria debe negociarse con el cliente; un sistema que registre asignaciones y contratos facilita esa gestión.
\end{itemize}

\section{Propuesta de solución}

Como respuesta a lo anterior, se propone un \textbf{sitio web de gestión} que centralice las operaciones claves de la empresa. Las funciones principales son:

\begin{enumerate}
  \item Asignar lugares y volumen de personal, a enviar, según tipo de cliente y necesidades.
  \item Asignar personal individualmente a cada ubicación, con instrucciones, herramientas y horarios específicos.
  \item Administrar contratos de personal de aseo, incluyendo salarios y ajustes del IPC.
  \item Administrar contratos y procesos de contratación de personal ejecutivo.
\end{enumerate}

En este trabajo de ingeniería de software, la solución que se se implementaría es una aplicación web con arquitectura cliente-servidor: un frontend React accesible desde navegador y un backend Express que persiste la información en PostgreSQL.

\section{Composición del Informe}

Este trabajo se organiza de la siguiente manera:

\begin{itemize}
  \item \textbf{Capítulo II} define los objetivos, requerimientos funcionales y no funcionales, límites e interfaces del software.
  \item \textbf{Capítulo III} presenta el análisis y diseño: modelo de datos, actores, casos de uso e interfaz de usuario.
  \item \textbf{Capítulo IV} describe la arquitectura y estructura del código desarrollado.
  \item \textbf{Capítulo V} expone las conclusiones del proyecto.
  \item \textbf{Capítulo VI} lista la bibliografía consultada.
  \item \textbf{Anexos} incluyen el instrumento de levantamiento y la documentación de pruebas de la API.
\end{itemize}

\screenshotsmall{dibujo.jpeg}{Ilustración contexto del trabajo de aseo}

% =============================================================================
\clearpage
\chapter{Definición de Software}
% =============================================================================

\section{Objetivo General del Software propuesto}

Desarrollar un sistema web, que centralice la gestión operativa y contractual de la empresa de aseos, permitiendo evaluar el volumen de personal requerido, asignar trabajadores a clientes con turnos y herramientas específicas, ademas administrar contratos de personal administrativo, ejecutivo y operativo.

\section{Objetivos Específicos del Software propuesto}

\begin{enumerate}
  \item Implementar un módulo de autenticación que restrinja el acceso a usuarios registrados.
  \item Permitir la evaluación y registro del volumen de personal recomendado por cliente y ubicación.
  \item Facilitar la asignación y distribución de personal con turnos, fechas de inicio y tareas detalladas.
  \item Gestionar contratos de trabajadores de aseo con duración, salario y ajuste IPC.
  \item Registrar procesos de contratación de personal ejecutivo (administrador, contador auditor, prevencionista de riesgo, químico, ventas).
  \item Ofrecer un dashboard de navegación hacia los módulos principales del sistema.
\end{enumerate}

\pagebreak
\section{Requerimientos Funcionales del Sw}

La Tabla~\ref{tab:rf} resume los requerimientos funcionales identificados.

\begin{longtable}{@{}p{0.9cm}p{3.9cm}p{2.3cm}p{2.8cm}p{3.8cm}@{}}
  \caption{Requerimientos funcionales del sistema} \label{tab:rf} \\
  \toprule
  \textbf{ID} & \textbf{Descripción} & \textbf{Actor} & \textbf{Ruta UI} & \textbf{Endpoint API} \\
  \midrule
  \endfirsthead
  \toprule
  \textbf{ID} & \textbf{Descripción} & \textbf{Actor} & \textbf{Ruta UI} & \textbf{Endpoint API} \\
  \midrule
  \endhead

  \small{RF1} & Autenticación de usuario mediante credenciales & Usuario administrativo & \scriptsize\texttt{/} & \scriptsize\texttt{POST /api/login} \\
  \small{RF2} & Evaluación de volumen de personal por cliente & Encargado de operaciones & \scriptsize\texttt{/asignacion} & \scriptsize\texttt{POST /api/crearAsignacion} \\
  \small{RF3} & Asignación de personal con turnos y tareas & Administrador de operaciones & \scriptsize\texttt{/asignacion/:id} & \scriptsize\texttt{POST /api/asignacion/:id/personal/agregar} \\
  \small{RF4} & Gestión de contratos de trabajadores & RRHH / Contador & \scriptsize\texttt{/contratos/personal} & \scriptsize\texttt{POST /api/contratos/personal} \\
  \small{RF5} & Contratación de personal ejecutivo & Gerente / Director & \scriptsize\texttt{/contratos/ejecutivos} & \scriptsize\texttt{POST /api/contratos/ejecutivos} \\
  \small{RF6} & Visualización de dashboard de módulos & Usuario autenticado & \scriptsize\texttt{/dashboard} & --- \\

  \bottomrule
\end{longtable}
\section{Límites}

El software propuesto contempla los siguientes límites de alcance:

\begin{itemize}
  \item Cobertura geográfica limitada a la Región del Bío-Bío; no incluye gestión de sucursales.
  \item No integra el portal de Mercado Público ni sistemas contables externos.
  \item No procesa pagos ni nómina; solo registra datos contractuales (salario, IPC).
  \item No envía notificaciones a trabajadores ni clientes.
  \item La edición de contratos ejecutivos en la interfaz es parcial (botón ``Modificar'' sin integración API).
  \item La validación del token JWT se realiza solo en el frontend; el backend no verifica el token en rutas protegidas.
\end{itemize}

\pagebreak
\section{Requerimientos No Funcionales del Sw}

\begin{table}[htbp]
  \centering
  \caption{Requerimientos no funcionales del sistema}
  \label{tab:rnf}
  \begin{tabular}{@{}p{2.5cm}p{10cm}@{}}
    \toprule
    \textbf{Categoría} & \textbf{Descripción} \\
    \midrule
    Seguridad & Contraseñas almacenadas con bcrypt; emisión de JWT en login; rate limit de 5 intentos fallidos cada 15 minutos en \texttt{/login}. \\
    Rendimiento & Respuestas JSON livianas; consultas directas a PostgreSQL sin capas intermedias innecesarias. \\
    Usabilidad & Interfaz en español, navegación lateral persistente, formularios con validación básica en cliente. \\
    Mantenibilidad & Separación en rutas, controladores y esquemas TypeORM; monorepo con \texttt{backend/} y \texttt{frontend/} independientes. \\
    Compatibilidad & Node.js $\geq$ 20, PostgreSQL $\geq$ 13, navegadores modernos con soporte ES6+. \\
    Disponibilidad & CI en GitHub Actions con verificación de sintaxis, build frontend y smoke test del backend. \\
    \bottomrule
  \end{tabular}
\end{table}

\pagebreak
\section{Interfaces externas de entrada}

\begin{table}[htbp]
  \centering
  \caption{Interfaces externas de entrada}
  \label{tab:entrada}
  \begin{tabular}{@{}p{3.5cm}p{9cm}@{}}
    \toprule
    \textbf{Interfaz} & \textbf{Descripción} \\
    \midrule
    Formulario de login & Usuario y contraseña enviados como JSON a \texttt{POST /api/login}. \\
    Formulario de asignación & Cliente, ubicación, necesidad y cantidad de personal recomendado. \\
    Formulario de distribución & ID de trabajador, turno, fecha de inicio, duración y detalles de tareas. \\
    Formulario de contratos personal & Nombre, fecha inicio, duración en años, salario e IPC. \\
    Formulario de contratos ejecutivo & Nombre, cargo, experiencia, salario, especialización, fecha entrevista, prioridad. \\
    Variables de entorno & \texttt{HOST}, \texttt{PORT}, \texttt{DATABASE}, \texttt{JWT\_SECRET}, credenciales DB. \\
    \bottomrule
  \end{tabular}
\end{table}

\section{Interfaces externas de Salida}

\begin{table}[htbp]
  \centering
  \caption{Interfaces externas de salida}
  \label{tab:salida}
  \begin{tabular}{@{}p{3.5cm}p{9cm}@{}}
    \toprule
    \textbf{Interfaz} & \textbf{Descripción} \\
    \midrule
    Respuestas JSON & Datos de asignaciones, contratos, distribución y autenticación en formato JSON. \\
    Vistas React & Páginas renderizadas: login, dashboard, asignación, distribución, contratos. \\
    Estados calculados & \texttt{Activo}, \texttt{Pendiente}, \texttt{Inactivo} (distribución); \texttt{Vencido}, \texttt{Próximo a vencer} (contratos). \\
    Herramientas sugeridas & Listado dinámico según ubicación (hospital $\rightarrow$ bioseguridad; otros $\rightarrow$ equipos estándar). \\
    Token JWT & Cadena firmada con vigencia de 8 horas, almacenada en \texttt{localStorage}. \\
    \bottomrule
  \end{tabular}
\end{table}

% =============================================================================
\clearpage
\chapter{Análisis y Diseño}
% =============================================================================

\section{Modelo de datos}

El sistema utiliza PostgreSQL como motor de base de datos, gestionado mediante TypeORM con esquemas \texttt{EntitySchema}. La Figura~\ref{fig:er} presenta el modelo entidad-relación lógico.

\begin{figure}[htbp]
  \centering
  \begin{tikzpicture}[
    entity/.style={draw, rectangle, minimum width=4.5cm, minimum height=1.5cm, 
                   align=center, font=\small, inner sep=5pt},
    rel/.style={-{Stealth}, thick}
  ]
    \node[entity] (usuarios) {usuarios\\\scriptsize id, username, password,\\nombre, rol, email};

    % asignaciones
    \node[entity, below left=2.5cm and 1.5cm of usuarios] (asignaciones) 
      {asignaciones\\\scriptsize id, cliente, ubicacion,\\necesidad, personal\_recomendado, estado};

    % contratos_p
    \node[entity, below=5.2cm of usuarios] (contratos_p) 
      {contratos\_personal\\\scriptsize id, nombre, inicio,\\duracion\_anos, salario, ipc};

    % asig_pers
    \node[entity, right=5.2cm of asignaciones] (asig_pers) 
      {asignaciones\_personal\\\scriptsize id, asignacion\_id,\\contrato\_personal\_id, turno,\\inicio, duracion, detalles};

    % contratos_e
    \node[entity, below=1.2cm of contratos_p] (contratos_e) 
      {contratos\_ejecutivo\\\scriptsize id, nombre, cargo, experiencia,\\salario, especializacion, estado};

    \draw[rel] (asignaciones) -- node[above, font=\scriptsize]{1:N} (asig_pers);
    \draw[rel] (contratos_p) -- node[above, font=\scriptsize]{1:N} (asig_pers);
  \end{tikzpicture}
  \caption{Modelo entidad-relación del sistema de gestión}
  \label{fig:er}
\end{figure}

Las relaciones entre \texttt{asignaciones\_personal} y las tablas padre son lógicas (sin claves foráneas declaradas en TypeORM). El campo \texttt{contrato\_personal\_id} vincula cada distribución con un trabajador contratado.
\pagebreak

\section{Actores Casos de uso}

\begin{table}[htbp]
  \centering
  \caption{Actores del sistema}
  \label{tab:actores}
  \begin{tabular}{@{}p{4cm}p{9cm}@{}}
    \toprule
    \textbf{Actor} & \textbf{Descripción} \\
    \midrule
    Usuario administrativo & Personal de oficina con acceso al sistema; inicia sesión y navega los módulos. \\
    Encargado de operaciones & Evalúa volumen de personal y crea asignaciones por cliente. \\
    Administrador RRHH & Gestiona contratos de trabajadores de aseo y sus condiciones. \\
    Gerente / Director & Inicia procesos de contratación de personal ejecutivo. \\
    \bottomrule
  \end{tabular}
\end{table}

\begin{table}[htbp]
  \centering
  \caption{Lista de casos de uso}
  \label{tab:cu}
  \begin{tabular}{@{}clp{8cm}@{}}
    \toprule
    \textbf{ID} & \textbf{Nombre} & \textbf{Descripción breve} \\
    \midrule
    CU1 & Iniciar sesión & Validar credenciales y obtener token JWT. \\
    CU2 & Evaluar volumen de personal & Registrar cliente, ubicación y cantidad recomendada. \\
    CU3 & Distribuir personal & Asignar trabajadores con turno, inicio y tareas. \\
    CU4 & Remover personal de asignación & Quitar un trabajador de la distribución activa. \\
    CU5 & Gestionar contratos personal & Crear y consultar contratos de trabajadores. \\
    CU6 & Iniciar contratación ejecutiva & Registrar candidato y proceso de entrevista. \\
    CU7 & Ver dashboard & Acceder al panel de módulos principales. \\
    CU8 & Cerrar sesión & Eliminar token y volver al login. \\
    \bottomrule
  \end{tabular}
\end{table}

\pagebreak
\section{Diagrama de casos de uso}

\begin{figure}[htbp]
  \centering
  \begin{tikzpicture}[
    actor/.style={draw, shape=rectangle, rounded corners, minimum height=0.9cm, font=\small},
    usecase/.style={draw, ellipse, minimum width=2.8cm, minimum height=0.8cm, align=center, font=\scriptsize},
    line/.style={-}
  ]
    \node[actor] (admin) at (-5, 3) {Usuario admin.};
    \node[actor] (ops) at (-5, 0) {Encargado ops.};
    \node[actor] (rrhh) at (-5, -3) {Admin. RRHH};

    \node[usecase] (cu1) at (0, 4) {CU1: Login};
    \node[usecase] (cu7) at (0, 2.5) {CU7: Dashboard};
    \node[usecase] (cu8) at (0, 1) {CU8: Cerrar sesión};
    \node[usecase] (cu2) at (0, -0.5) {CU2: Evaluar volumen};
    \node[usecase] (cu3) at (0, -2) {CU3: Distribuir personal};
    \node[usecase] (cu4) at (-5, -2) {CU4: Remover personal};
    \node[usecase] (cu5) at (5.5, -2.5) {CU5: Contratos personal};
    \node[usecase] (cu6) at (5.5, -3.5) {CU6: Contratación ejecutiva};

    \foreach \a/\cus in {admin/{cu1,cu7,cu8}, ops/{cu2,cu3,cu4}, rrhh/{cu5,cu6}} {
      \foreach \cu in \cus {
        \draw[line] (\a) -- (\cu);
      }
    }
  \end{tikzpicture}
  \caption{Diagrama general de casos de uso}
  \label{fig:cu}
\end{figure}

\pagebreak
\section{Especificación de los Casos de Uso}

\subsection{CU1 -- Iniciar sesión}

\begin{table}[htbp]
  \centering
  \caption{Especificación CU1: Iniciar sesión}
  \begin{tabular}{@{}p{3.5cm}p{10cm}@{}}
    \toprule
    \textbf{Campo} & \textbf{Descripción} \\
    \midrule
    ¿Quién? & Usuario del sistema (personal administrativo o ejecutivo). \\
    ¿Qué? & Iniciar sesión en el sistema de gestión. \\
    ¿Cuándo? & Al acceder al sistema desde la página principal (\texttt{/}). \\
    ¿Cómo? & Ingresar credenciales válidas en el formulario de login. \\
    Restricción & Solo usuarios registrados; máximo 5 intentos fallidos por ventana de 15 minutos. \\
    Postcondición & Usuario autenticado accede al dashboard con token JWT almacenado. \\
    \bottomrule
  \end{tabular}
\end{table}

\subsection{CU2 -- Evaluar volumen de personal}

\begin{table}[htbp]
  \centering
  \caption{Especificación CU2: Evaluar volumen de personal}
  \begin{tabular}{@{}p{3.5cm}p{10cm}@{}}
    \toprule
    \textbf{Campo} & \textbf{Descripción} \\
    \midrule
    ¿Quién? & Administrador o encargado de operaciones. \\
    ¿Qué? & Evaluar y registrar el volumen de personal necesario para un cliente. \\
    ¿Cuándo? & Al recibir una nueva solicitud de servicio o cambio en necesidades. \\
    ¿Cómo? & Seleccionar tipo de cliente, ubicación, describir necesidad y cantidad recomendada. \\
    Restricción & La evaluación debe considerar el tipo de cliente (hospital requiere más personal que banco). \\
    Postcondición & Evaluación guardada y disponible para asignación de personal. \\
    \bottomrule
  \end{tabular}
\end{table}

\pagebreak
\subsection{CU3 -- Distribuir personal}

\begin{table}[htbp]
  \centering
  \caption{Especificación CU3: Asignar personal a cliente}
  \begin{tabular}{@{}p{3.5cm}p{10cm}@{}}
    \toprule
    \textbf{Campo} & \textbf{Descripción} \\
    \midrule
    ¿Quién? & Administrador de operaciones. \\
    ¿Qué? & Asignar personal específico con equipamiento y turnos. \\
    ¿Cuándo? & Después de la evaluación de volumen o ante cambios en distribución. \\
    ¿Cómo? & Seleccionar trabajador, turno, fecha de inicio, duración y tareas específicas. \\
    Restricción & Solo trabajadores con contrato activo; equipamiento según tipo de ubicación. \\
    Postcondición & Asignación registrada en \texttt{asignaciones\_personal}. \\
    \bottomrule
  \end{tabular}
\end{table}

\subsection{CU5 -- Gestionar contratos de personal}

\begin{table}[htbp]
  \centering
  \caption{Especificación CU5: Gestión de contratos de trabajo}
  \begin{tabular}{@{}p{3.5cm}p{10cm}@{}}
    \toprule
    \textbf{Campo} & \textbf{Descripción} \\
    \midrule
    ¿Quién? & Administrador de RRHH o contador. \\
    ¿Qué? & Crear y consultar contratos de trabajo para trabajadores. \\
    ¿Cuándo? & Al contratar nuevo personal o renovar contratos existentes. \\
    ¿Cómo? & Ingresar nombre, fecha de inicio, duración (1--5 años), salario e IPC. \\
    Restricción & Contratos largos deben incluir ajuste IPC; duración máxima 5 años. \\
    Postcondición & Contrato creado y almacenado en base de datos. \\
    \bottomrule
  \end{tabular}
\end{table}

\pagebreak
\section{Diseño interfaz y navegación}

El flujo de navegación del sistema sigue la estructura definida en \texttt{App.jsx}:

\begin{center}
  \texttt{/} $\rightarrow$ Login $\rightarrow$ \texttt{/dashboard} $\rightarrow$ \texttt{/asignacion} $\rightarrow$ \texttt{/asignacion/:id}\\
  \texttt{/dashboard} $\rightarrow$ \texttt{/contratos/personal} \quad | \quad \texttt{/contratos/ejecutivos}
\end{center}

El menú lateral (\texttt{GeneracionPagina.jsx}) permanece visible en todas las páginas autenticadas, con resaltado del módulo activo.

\subsection{Guías de estilos}

La interfaz adopta un patrón de \textbf{panel de administración} con tres zonas: cabecera superior fija, barra lateral de navegación y área de contenido principal. Los formularios utilizan etiquetas sobre campos de entrada, botones de acción primaria en color sólido y tablas para listados de datos. El diseño es responsivo mediante clases de Tailwind CSS (\texttt{md:}, \texttt{flex-col}).

\subsection{Guía de colores}

\begin{table}[htbp]
  \centering
  \caption{Paleta de colores de la interfaz}
  \begin{tabular}{@{}p{3cm}p{3cm}p{7cm}@{}}
    \toprule
    \textbf{Elemento} & \textbf{Clase Tailwind} & \textbf{Uso} \\
    \midrule
    Primario & \texttt{bg-blue-600} & Cabecera, botones principales, ítem activo del menú. \\
    Fondo & \texttt{bg-gray-100} & Fondo general de las páginas. \\
    Tarjetas & \texttt{bg-white} & Contenedores de formularios y módulos del dashboard. \\
    Texto secundario & \texttt{text-gray-500} & Descripciones y etiquetas de navegación. \\
    Estados activos & \texttt{bg-green-100}, \texttt{bg-yellow-100} & Badges de estado en tablas de contratos y distribución. \\
    \bottomrule
  \end{tabular}
\end{table}

\subsection{Tipografía}

Se utiliza la fuente sans-serif por defecto del sistema operativo, aplicada mediante Tailwind CSS. Los títulos de página emplean \texttt{text-xl font-bold}; las etiquetas de formulario usan \texttt{text-sm font-medium}; la navegación lateral muestra el nombre del módulo en \texttt{font-medium} y una descripción en \texttt{text-xs}.

\subsection{Composición de las interfaces}

\screenshot{Screenshot from 2026-07-02 22-08-48.png}{Interfaz de inicio de sesión del sistema}

\screenshot{Screenshot from 2026-07-02 22-09-18.png}{Dashboard de gestión con acceso a los módulos principales}

\screenshot{Screenshot from 2026-07-02 22-09-00.png}{Página de evaluación de volumen de personal y listado de asignaciones}

\screenshot{Screenshot from 2026-07-02 22-09-03.png}{Página de distribución de personal para una asignación específica}

\screenshot{Screenshot from 2026-07-02 22-09-05.png}{Gestión de contratos de trabajadores de aseo}

\screenshot{Screenshot from 2026-07-02 22-09-07.png}{Módulo de contratación de personal ejecutivo}

\screenshot{Screenshot from 2026-07-02 22-08-53.png}{Vista complementaria del módulo de asignaciones}

\screenshot{Screenshot from 2026-07-02 22-08-57.png}{Vista complementaria de la distribución de personal}

% =============================================================================
\clearpage
\chapter{Desarrollo}
% =============================================================================

\section{Diseño de arquitectura}

La aplicación sigue una arquitectura de tres capas desacopladas. La Figura~\ref{fig:arquitectura} ilustra el flujo de comunicación entre componentes.

\begin{figure}[htbp]
  \centering
  \begin{tikzpicture}[
    box/.style={draw, rectangle, rounded corners, minimum width=2.8cm, minimum height=1cm, align=center, font=\small},
    db/.style={draw, cylinder, shape border rotate=90, aspect=0.25, minimum height=1.2cm, minimum width=1.5cm, align=center, font=\small},
    arr/.style={-{Stealth}, thick}
  ]
    \node[box, fill=blue!10] (fe) at (0, 0) {Frontend\\React + Vite\\:5173};
    \node[box, fill=green!10] (be) at (5, 0) {Backend\\Express + TypeORM\\:3000};
    \node[db, fill=orange!10] (db) at (10, 0) {(PostgreSQL)};

    \draw[arr] (fe) -- node[above, font=\scriptsize]{HTTP/JSON + JWT} (be);
    \draw[arr] (be) -- node[above, font=\scriptsize]{SQL} (db);
  \end{tikzpicture}
  \caption{Arquitectura general del sistema}
  \label{fig:arquitectura}
\end{figure}

\textbf{Patrones y decisiones de diseño:}
\begin{itemize}
  \item \textbf{Monorepo funcional:} Directorios \texttt{backend/} y \texttt{frontend/} con dependencias independientes gestionadas por pnpm.
  \item \textbf{Capas en backend:} Rutas $\rightarrow$ Controladores $\rightarrow$ Repositorios TypeORM.
  \item \textbf{API REST:} Endpoints JSON bajo el prefijo \texttt{/api}.
  \item \textbf{SPA:} El frontend es una Single Page Application con enrutamiento del lado del cliente.
  \item \textbf{Seed de datos:} Al iniciar el servidor se cargan datos de prueba si las tablas están vacías.
  \item \textbf{CI/CD:} GitHub Actions ejecuta lint, build y smoke test en cada push a \texttt{master}.
\end{itemize}

\pagebreak
\section{Estructura del código}

\subsection{Backend}

\begin{lstlisting}[language=bash, caption={Estructura del directorio backend/src/}]
backend/src/
  servidor.js              # Bootstrap Express, CORS, montaje de routers
  configuracion/env.js     # Variables de entorno
  database/database.js     # Conexion TypeORM y seed de datos
  schema/                  # 5 EntitySchema (usuarios, asignaciones, etc.)
  routes/                  # login, asignaciones, contratospersonal, contratosejecutivo
  controllers/             # Logica de negocio y estados calculados
\end{lstlisting}

\textbf{Archivos clave:}
\begin{itemize}
  \item \texttt{servidor.js}: Inicializa Express, habilita CORS para desarrollo, monta los cuatro routers bajo \texttt{/api} y conecta PostgreSQL antes de escuchar en el puerto configurado.
  \item \texttt{asignaciones.controller.js}: Implementa \texttt{herramientasDeAcuerdoAUbicacion()} (bioseguridad para hospitales) y \texttt{calcularEstado()} para distribuciones según fecha de inicio y duración.
  \item \texttt{contratospersonal.controller.js}: Calcula estados contractuales (\texttt{Pendiente}, \texttt{Activo}, \texttt{Próximo a vencer}, \texttt{Vencido}).
  \item \texttt{login.controller.js}: Valida credenciales con bcrypt y emite JWT con vigencia de 8 horas.
\end{itemize}

\pagebreak
\subsection{Frontend}

\begin{lstlisting}[language=bash, caption={Estructura del directorio frontend/src/}]
frontend/src/
  main.jsx                 # Punto de entrada React
  App.jsx                  # Rutas y guard de autenticacion
  api.js                   # Cliente REST centralizado
  Login.jsx                # Pagina de inicio de sesion
  Dashboard.jsx            # Panel de modulos
  Asignacion.jsx           # Evaluacion de volumen
  DistribucionAsignacion.jsx  # Distribucion de personal
  ContratosPersonal.jsx    # Contratos de trabajadores
  ContratosEjecutivo.jsx   # Contratacion ejecutiva
  GeneracionPagina.jsx     # Header, Sidebar, Layout compartido
  style.css                # Import de Tailwind CSS
\end{lstlisting}

\textbf{Stack frontend:} React 19, Vite 8, react-router-dom 7, Tailwind CSS 4.

El archivo \texttt{api.js} centraliza todas las llamadas HTTP con la función \texttt{requestAPIinterna()}, que adjunta automáticamente el header \texttt{Authorization: Bearer <token>} cuando existe sesión activa.

% =============================================================================
\clearpage
\chapter{Conclusión del proyecto}
% =============================================================================

El proyecto entrega un prototipo funcional de sistema de gestión para una empresa de aseos de la Región del Bío-Bío. Se implementaron con éxito los módulos de autenticación, evaluación de volumen de personal, distribución de trabajadores, gestión de contratos de personal y registro de contratación ejecutiva, con persistencia en PostgreSQL y una interfaz web operativa.

Los principales logros del desarrollo son:
\begin{itemize}
  \item API REST con 14 endpoints documentados y probados.
  \item Modelo de datos con cinco entidades que cubren usuarios, asignaciones, distribución y contratos.
  \item Interfaz de usuario coherente con navegación lateral y capturas de pantalla que evidencian su funcionamiento.
  \item Medidas de seguridad básicas: hash de contraseñas, JWT y rate limiting en login.
  \item Pipeline de integración continua con verificación automática de sintaxis, build y disponibilidad del backend.
\end{itemize}

Como limitaciones pendientes se identifican: la ausencia de middleware de validación JWT en el backend, la falta de tests unitarios automatizados, la edición incompleta de contratos ejecutivos en la UI, y la configuración de CORS abierta (\texttt{origin: '*'}) orientada solo a desarrollo.

Como trabajo futuro se recomienda: implementar control de acceso basado en roles, desplegar en un entorno de producción con HTTPS, integrar notificaciones a trabajadores, y conectar con sistemas contables o el portal de Mercado Público según las necesidades de la empresa.

% =============================================================================
\clearpage
\chapter{Bibliografía}
% =============================================================================

\begin{enumerate}
  \item Meta Open Source. (2024). \textit{React Documentation}. \url{https://react.dev/}
  \item Express.js. (2024). \textit{Express -- Node.js web application framework}. \url{https://expressjs.com/}
  \item TypeORM. (2024). \textit{TypeORM Documentation}. \url{https://typeorm.io/}
  \item PostgreSQL Global Development Group. (2024). \textit{PostgreSQL Documentation}. \url{https://www.postgresql.org/docs/}
  \item Tailwind Labs. (2024). \textit{Tailwind CSS Documentation}. \url{https://tailwindcss.com/docs}
  \item IETF. (2015). \textit{RFC 7519: JSON Web Token (JWT)}. \url{https://datatracker.ietf.org/doc/html/rfc7519}
  \item Vite Team. (2024). \textit{Vite Guide}. \url{https://vite.dev/guide/}
  \item ChileCompra. (2024). \textit{Portal Mercado Público}. \url{https://www.mercadopublico.cl/}
  \item Instituto Nacional de Estadísticas de Chile. (2024). \textit{Índice de Precios al Consumidor (IPC)}. \url{https://www.ine.cl/}
  \item Sommerville, I. (2015). \textit{Software Engineering} (10th ed.). Pearson Education.
\end{enumerate}

% =============================================================================
\appendix
\chapter{Instrumento y datos recopilados de análisis del contexto}
% =============================================================================

\section{Descripción del contexto organizacional}

\textit{Extracto de levantamiento de información:}

La empresa de aseos atiende clientes de la Región del Bío-Bío, ofreciendo servicios para instituciones públicas (hospitales, bancos, universidades) a través del mercado público y clientes privados por comunicación directa. Se provee personal de aseo con herramientas básicas; la cantidad de trabajadores depende del tipo de cliente y sus necesidades específicas.

Los contratos varían entre 1 y 5 años. La empresa mantiene clientes frecuentes no exclusivos, pudiendo incorporar personal ante nuevas construcciones o redistribuir trabajadores en cambios menores. La oficina está en Concepción y no hay proyección de expansión a otras regiones.

El personal fijo incluye dos contadores, un encargado de ventas y personal de operaciones. Cada cliente requiere un costeo distinto y las desviaciones presupuestarias deben negociarse, considerando ajustes por IPC en contratos plurianuales.

\section{Propuesta inicial}

\textit{Extracto de propuesta de solución:}

\begin{itemize}
  \item Asignar lugares y volumen de personal a enviar.
  \item Asignar personal individualmente para una ubicación con instrucciones, herramientas y horamientos específicos.
  \item Administrar contratos de personal y pagos.
  \item Administrar contratos ejecutivos.
\end{itemize}

% =============================================================================
\clearpage
\chapter{Pruebas de API}
% =============================================================================

\section{Configuración del entorno de pruebas}

\textbf{Credenciales de prueba:} usuario \texttt{admin}, contraseña \texttt{password}.

\textbf{Ejecución del backend (Linux):}
\begin{lstlisting}[language=bash]
bash prueba_backend.sh
\end{lstlisting}

\textbf{Ejecución del frontend (Linux):}
\begin{lstlisting}[language=bash]
bash prueba_frontend.sh
\end{lstlisting}

El script \texttt{prueba\_backend.sh} levanta un contenedor Docker con PostgreSQL 13, configura las variables de entorno y ejecuta el servidor en modo desarrollo con nodemon.

\section{Endpoints de autenticación}

\begin{longtable}{@{}p{1.5cm}p{4cm}p{8cm}@{}}
  \caption{Endpoint de autenticación} \\
  \toprule
  \textbf{Metodo} & \textbf{Ruta} & \textbf{Descripcion} \\
  \midrule
  \endfirsthead
  POST & \texttt{/api/login} & Valida credenciales y retorna JWT. \\
  \bottomrule
\end{longtable}

\textbf{Ejemplo curl:}
\begin{lstlisting}[language=bash]
curl -X POST http://localhost:3000/api/login \
  -H "Content-Type: application/json" \
  -d '{"usuario":"admin","contraseña":"password"}'
\end{lstlisting}

\textbf{Respuesta esperada:}
\begin{lstlisting}
{
  "success": true,
  "token": "eyJhbGciOiJIUzI1NiIs...",
  "user": { "id": 1, "nombre": "Administrador", "rol": "admin" }
}
\end{lstlisting}

\section{Endpoints de asignaciones}

\begin{longtable}{@{}p{1.5cm}p{5.5cm}p{6.5cm}@{}}
  \caption{Endpoints del módulo de asignaciones} \\
  \toprule
  \textbf{Metodo} & \textbf{Ruta} & \textbf{Descripcion} \\
  \midrule
  \endfirsthead
  GET & \scriptsize\texttt{/api/asignacion} & Lista todas las asignaciones con resumen de personal. \\
  POST & \scriptsize\texttt{/api/crearAsignacion} & Crea una nueva asignación. \\
  GET & \scriptsize\texttt{/api/asignacion/:id/detalles} & Detalles, herramientas y estado. \\
  GET & \scriptsize\texttt{/api/asignacion/:id/distribucion} & Distribución de personal asignado. \\
  POST & \scriptsize\texttt{/api/asignacion/:id/personal/agregar} & Agrega trabajador a asignación. \\
  POST & \scriptsize\texttt{/api/asignacion/:id/personal/remover} & Remueve trabajador de asignación. \\
  GET & \scriptsize\texttt{/api/asignacion/:id/personal/detalles} & Detalles de tareas de un trabajador. \\
  \bottomrule
\end{longtable}

\textbf{Ejemplo -- crear asignación:}
\begin{lstlisting}[language=bash]
curl -X POST http://localhost:3000/api/crearAsignacion \
  -H "Content-Type: application/json" \
  -d '{
    "cliente": "Hospital",
    "ubicacion": "Concepcion - Pabellon Central",
    "necesidad": "Limpieza diaria de areas comunes",
    "personal": 4
  }'
\end{lstlisting}

\textbf{Ejemplo -- listar asignaciones:}
\begin{lstlisting}[language=bash]
curl http://localhost:3000/api/asignacion
\end{lstlisting}

\pagebreak
\section{Endpoints de contratos de personal}

\begin{longtable}{@{}p{1.5cm}p{4.5cm}p{6.5cm}@{}}
  \caption{Endpoints del módulo de contratos de personal} \\
  \toprule
  \textbf{Metodo} & \textbf{Ruta} & \textbf{Descripcion} \\
  \midrule
  \endfirsthead
  GET & \scriptsize\texttt{/api/contratos/personal} & Lista IDs y nombres de trabajadores. \\
  GET & \scriptsize\texttt{/api/contratos/personal/:id} & Detalle completo con estado calculado. \\
  POST & \scriptsize\texttt{/api/contratos/personal} & Crea un nuevo contrato. \\
  POST & \scriptsize\texttt{/api/contratos/personal/:id} & Modifica un contrato existente. \\
  \bottomrule
\end{longtable}

\textbf{Ejemplo -- crear contrato:}
\begin{lstlisting}[language=bash]
curl -X POST http://localhost:3000/api/contratos/personal \
  -H "Content-Type: application/json" \
  -d '{
    "nombre": "Pedro Soto",
    "inicio": "2026-01-01",
    "duracion": 3,
    "salario": 380000,
    "ipc": 3.5
  }'
\end{lstlisting}

\section{Endpoints de contratos ejecutivos}

\begin{longtable}{@{}p{1.5cm}p{4cm}p{6.5cm}@{}}
  \caption{Endpoints del módulo de contratos ejecutivos} \\
  \toprule
  \textbf{Metodo} & \textbf{Ruta} & \textbf{Descripcion} \\
  \midrule
  \endfirsthead
  GET & \scriptsize\texttt{/api/contratos/ejecutivos} & Lista todos los registros ejecutivos. \\
  POST & \scriptsize\texttt{/api/contratos/ejecutivos} & Inicia proceso de contratación. \\
  \bottomrule
\end{longtable}

\textbf{Ejemplo -- crear registro ejecutivo:}
\begin{lstlisting}[language=bash]
curl -X POST http://localhost:3000/api/contratos/ejecutivos \
  -H "Content-Type: application/json" \
  -d '{
    "nombre": "Laura Fernandez",
    "cargo": "Contador",
    "experiencia": 8,
    "salario": 1000000,
    "especializacion": "Tributaria",
    "fecha_entrevista": "2026-08-15",
    "prioridad": "Alta"
  }'
\end{lstlisting}

\section{Resumen de endpoints}

En total el sistema expone \textbf{14 endpoints} bajo el prefijo \texttt{/api}, distribuidos en cuatro módulos: autenticación (1), asignaciones (7), contratos de personal (4) y contratos ejecutivos (2).

\end{document}